\paragraph{First-order equivariant class about flat space: the pairing of the action, commutant labels and the double zeros (conditional; not closure).}
Take the flat point of the preceding class, \(\hat X_a=\mathrm iu_a{}^b\partial_b\) on \(L^2(\mathbb R^4\times\mathrm{SO}(3,1))\) with the reference measure, \(C_{\mu\nu\lambda}=0\), no matter term, at quadratic order in fluctuations. With \(p_b\) the eigenvalue of \(-\mathrm i\partial_b\), \(\hat X_a\) is multiplication by \(-u_a{}^bp_b\), and \((p,u)\mapsto(\tilde p=up,u)\) has unit Jacobian, so \(\hat X_a\) is unitarily equivalent to \(-\tilde p_a\otimes1\) on \(L^2(\mathbb R^4)\otimes L^2(\mathrm{SO}(3,1))\). Uniformly distributed \(p_{(a)}\) of this kind appear for \(\mathrm{Spin}(n)\) in~\cite{Hanada2006}, which then applies the quenched reduced model. At this point the matrices have joint spectrum \(\mathbb R^4\) with uniform multiplicity: apart from that multiplicity they do not distinguish the class from a commuting background without a fibre, and the algebra \(\mathcal A\) remains data. For Hermitian Hilbert--Schmidt fluctuations \(A_a\), with kernels between points \(m=(p,u)\) and \(m'=(p',u')\), let \(Q_a(m,m')=-u_a{}^bp_b+u'_a{}^bp'_b\), which is real. Equation~\eqref{eq:Smatrix} gives
\[
S_2=\frac1{2g_{\mathrm{YM}}^2}\,\eta^{ab}\mathrm{Tr}\bigl(A_a(\mathcal EA)_b\bigr),\qquad
(\mathcal EA)_b(m,m')=Q\cdot Q\,A_b(m,m')-Q_b\,Q^aA_a(m,m'),
\]
and the same form holds exactly for commuting \(N\times N\) matrices, with \(Q_a\) the differences of eigenvalues. So \(\mathcal EA=0\) at non-null \(Q\) forces \(A_a\propto Q_a\), which is the gauge direction; at null \(Q\neq0\) it forces \(Q\cdot A=0\), and the form \(\eta^{ab}\bar A_aA_b\) is positive definite on \(Q^\perp/Q\) at each pair \((m,m')\). Modulo the gauge directions \(A_a=Q_a\Lambda\), on which the form has the value \(Q\cdot Q\,|\Lambda|^2\) of either sign, the form is indefinite on these kernels only off shell and at \(Q=0\), the moduli of the commuting set. This statement concerns the signature, that is the relative sign of modes; the overall sign, relative to matter and to the weight of the integral, is not addressed. In background gauge the operator is multiplication by \(Q\cdot Q\): on Hilbert--Schmidt kernels it is self-adjoint, with no eigenvector and no Jordan block, and every probe with polarisation of positive norm has a positive spectral measure. At finite \(N\) this pairing is the trace form in which the action is written, so no regulator enters there; the truncations of first-order operators below are a choice. For a commuting background without a fibre, positivity in the bilocal description is stated in~\cite{HoKawaiSteinacker2025} for the planar sector with quenched diagonal elements, and a no-ghost theorem based on the same kind of inner product is given in~\cite{Steinacker2019NoGhost} on a different, noncommutative background.

The first-order operators sit on this space as singular kernels at \(u=u'\), and they can be labelled. The operators commuting with all \(\hat X_a\) include the functions of \(u\), the \(\mathcal P_g=-\mathrm i\partial_g\) and, for \(\lambda\in\mathfrak{so}(3,1)\), the derivation \(\mathcal J_\lambda=\tfrac12\lambda^{cd}u^{r}{}_cu^{s}{}_d\mathcal S_{rs}-y_g\lambda^{gc}\partial_c\), which generates a unitary right translation of \(u\) combined with a Lorentz transformation of \(y\). The last two are the right-invariant fields on \(Z=\mathrm{ISO}(3,1)\); the identification of the flat bundle with the Poincar\'e group, and the commuting right action on a bundle that is a Lie group, are in~\cite{Hanada2006}. If \(\mathcal C\) commutes with all \(\hat X_a\), then the linearised operator obeys \(\mathcal E_b[\tfrac12\{f,\mathcal C\}]=\tfrac12\{\mathcal E_b[f],\mathcal C\}\) for every \(f\), since it is built from \(\mathrm{ad}\,\hat X_a\) and \(\mathrm{ad}\,[\hat X_a,\hat X_b]\). In section values, write
\[
h_b{}^c=f_b{}^c+y_gf_b{}^{gc},\qquad\omega_b{}^{cd}=f_b{}^{cd},
\]
with arbitrary functions \(f_b{}^c\), \(f_b{}^{cd}=-f_b{}^{dc}\) of \(y\), and let \(F\) denote \(\partial_af_b-\partial_bf_a\) on the first index of each. Then, exactly and for arbitrary functions,
\[
\mathsf T_{ab}{}^c=F_{ab}{}^c+y_gF_{ab}{}^{gc},\qquad\mathsf R_{ab}{}^{cd}=F_{ab}{}^{cd},
\]
and with \(\mathcal E^{(1)}{}_b{}^e=-\partial^a\mathsf T_{ab}{}^e+\mathsf R_{ab}{}^{ae}\), \(\mathcal E^{(2)}{}_b{}^{de}=\partial^a\mathsf R_{ab}{}^{de}\), \(\mathcal E^{(3)}{}_b=\partial^a(\partial_a\mathsf t_b-\partial_b\mathsf t_a)\), the linearised brackets of the matrix equation,
\[
\mathcal E^{(2)}{}_b{}^{de}=\partial^aF_{ab}{}^{de},\qquad
\mathcal E^{(1)}{}_b{}^e+y_g\mathcal E^{(2)}{}_b{}^{ge}=-\partial^aF_{ab}{}^e,\qquad
2\mathcal E^{(3)}{}_b=\partial_c\mathcal E^{(1)}{}_b{}^c-\eta_{bg}\mathcal E^{(2)}{}_c{}^{gc}.
\]
So the forty equations about flat space are ten Maxwell equations, one for each label in \(\mathfrak{iso}(3,1)\), with \(A_a\) an eleventh, and with one gauge parameter each, the count \(6+4+1\) of the class; each label has two polarisations, so \(22=20+2\). Here \(\mathsf T\) and \(\mathsf R\) are the linearised torsion and curvature, which in the gauge theory of the Poincar\'e group are read as the translational and Lorentz field strengths~\cite{Hehl2017}. For fixed label all zeros at \(q^2=0\) are simple: \(\mathcal E_b[e^{\mathrm iq\cdot y}\mathcal C]\) with \(\mathcal C=\mathcal P_g\) or \(\mathcal J_\lambda\) is \(q^2\) times the mode on polarisations transverse to \(q\). The double zeros of the preceding paragraph are a statement at fixed eigenvalue of the translations and they stand as such; they are the nilpotent action \([\mathcal P_g,\mathcal J_\lambda]\propto\lambda_g{}^c\mathcal P_c\) of translations on the labels, whose rank at null \(q\) is that of \(\lambda\mapsto q_g\lambda^{gc}\), three, for each of the two polarisations. The redefinition contains \(y\) explicitly, so it is not translation invariant and it does not remove the solutions linear in \(y\). Those solutions are \(e^{\mathrm ik\cdot y}\mathcal J_\lambda\) at null \(k\), the derivative at \(t=0\) of \(e^{\mathrm ik\cdot y}\exp(t\mathcal J_\lambda)\), and these are unitary operators on the Hilbert space that, times a transverse polarisation \(u_a{}^r\varepsilon_r\), solve the same equation for all \(t\). The argument carries over where \(Z\) is a group with left-invariant background: on \(\mathrm{SO}(4,1)\) with the bilinear term the labels are \(\mathfrak{so}(4,1)\) and the identity \(\mathbb 1\), and every equivariant first-order fluctuation is a sum of terms \(\tfrac12\{f_a,\mathcal C\}\) with \(f_a\) a dressed one-form, so the same order-zero equation \(\nabla^a\mathsf F_{ab}=0\) of the preceding paragraph governs each label on the unit de Sitter space.

Each label retains its gauge parameter, so these modes stay massless at this order; that the fields of the expansion in differential operators are massless at tree level is stated in~\cite{Sakai2019}. The first two brackets of the matrix equation contain neither \(\mathsf F_{ab}\) nor order-zero fields, so the constant \(\Theta_{ab}\) of the preceding paragraph and singlets of a fixed factor do not change the linear equations of \(h\) and \(\omega\). The torsion modes are therefore not removed by anything found here. That is a bounded statement: quadratic order, background torsion zero, operators of the stated order and the commutant labels above.

The pairing of a label is not contracted with \(\eta\). A kernel \(A_a(m,m')=u_a{}^r\varepsilon_r\mathcal O(m,m')\), with \(u\) the fibre point of \(m\) and \(\eta^{rs}\varepsilon_r\varepsilon_s=1\), has pairing \(\eta^{ab}\mathrm{Tr}(A_a^\dagger A_b)=\int|\mathcal O|^2\), whatever \(\mathcal O\) is, and if \(\mathcal O\) carries nonzero translation charge then \(\eta^{ab}\mathrm{Tr}(A_aA_b)=0\), so the Hermitian part of \(A_a\) has half of that value. For \(\mathcal O\) a truncation of \(\tfrac12\{f,\mathcal P_g\}\) with a positive weight \(w(p)\) of unit total mass and finite second moments, and \(f\) slowly varying on the scale of \(w\), the pairing of the truncations for \(g\) and \(h\) is \(\int|f|^2\,\mathsf G_{gh}\) with \(\mathsf G_{gh}=\int w\,p_gp_h\), up to a positive fibre factor common to all \(g,h\) and infrared volumes. The matrix \(\mathsf G\) is positive semi-definite, positive definite unless \(w\) lies on a hyperplane through the origin, and a Lorentz transformation of the weight changes it by a congruence. It is never a nonzero multiple of \(\eta_{gh}\), since \(\mathsf G_{00}\ge0\) and \(\mathsf G_{11}\ge0\) while \(\eta_{00}\) and \(\eta_{11}\) have opposite signs. So in this pairing \(\omega_i{}^{0j}\), with \(i\) transverse in this sense and \(\mathcal O\) a truncation of \(f\mathcal J_\lambda\), has a positive value, the boost part of \(\omega_i\) included, and for \(w\) not on such a hyperplane so do \(h_i{}^0\) and the direction \(h_i{}^\mu\propto q^\mu\), whose value is not zero. This differs from the reading in~\cite{FurutaHanadaKawaiKimura2007}, where those components are expected to be negative were the equation derived from a field theory action, and from the statements in~\cite{HoKawaiSteinacker2025} on tensor fields, that the time component of each tensor field has negative norm and that, in a minimal model without a fibre, the transverse zero-norm components decouple; that reference records that this is very different from the positivity of bilocal fields and that there is no clear correspondence between the bilocal and the tensor inner products. The pairing here is of the bilocal kind, so this is a tension between readings and not a contradiction, and the truncation is one regulator-dependent map from that pairing to a form on the label coefficients, not a derivation of the tensor inner product. It does not settle the question and does not show that negative norm is absent for effective fields: the form diverges as the cutoff is removed, it depends on the frame and on \(w\), it is not an invariant form on \(\mathfrak{iso}(3,1)\), which has no positive invariant form, and the kinetic term of an effective field, such as one induced at one loop, need not equal it. The published caution therefore stands for the covariant reading.

The double zero itself can be followed. Fix a null \(k\), a one-parameter subgroup \(\Lambda_t=\exp(t\lambda)\) with \(k\cdot\lambda\neq0\) and \(\phi^2=e^{-t^2/\sigma^2}/(\sigma\sqrt\pi)\), and take the kernels \(e^{\mathrm ik\cdot y}\,w(\mathcal P)^{1/2}\int\mathrm dt\,\psi(t)R_{\Lambda_t}\), with \(R\) the right translation of \(u\) alone, not \(\exp(t\mathcal J_\lambda)\), on which the background-gauge operator is multiplication by \((k+(1-\Lambda_t)p)^2\). Let \(v^g\) be defined by \(\partial_t(k+(1-\Lambda_t)p)^2|_{t=0}=v^gp_g\), which is \(\pm2(k\cdot\lambda)^g\). The probes with \(\psi=\phi\), multiplied by \(v\cdot\mathcal P\) and normalised, and with \(\psi=\sqrt2\sigma\phi'\), are orthonormal, and they are cutoffs of the frame and fibre parts. The compression of the operator to their span is
\[
\begin{pmatrix}O(\sigma^2)&\varkappa\\\varkappa&O(\sigma^2)\end{pmatrix},\qquad
\varkappa=-\frac{\sigma}{\sqrt2}\sqrt{v^g\mathsf G_{gh}v^h}\,\bigl(1+O(\sigma^2)\bigr),
\]
as \(\sigma\to0\) at fixed \(w\), so its eigenvalues are real, with orthonormal eigenvectors: the resolvent of this compression has two simple poles with residues of one sign, which merge as \(\sigma\to0\). These are not poles of the two-point function of the probes, whose spectral measure under the full multiplication operator is positive and continuous. The entries of this compression are evaluated exactly on a Gaussian \(w\) for one boost and for one rotation on its covering line (\(\varkappa=-\sqrt2\sigma P\,(1+O(\sigma^2))\) for the Gaussian \(w\) of variance \(P^2\) in each component, \(v\) of Euclidean norm \(2\)), with a seeded Monte Carlo agreeing; smearing over the other five fibre directions is not treated. A Jordan block is read only in the basis where the second probe is divided by \(\sqrt2\sigma\), whose squared norm \(1/(2\sigma^2)\) diverges. The limits do not commute: the spectral width is of order \(\sigma|k|P+\sigma^2P^2\), and removing the momentum cutoff first spreads the measure.

None of this extends automatically. On any background \((\mathcal MA)_b=[\hat X^a,[\hat X_a,A_b]]-2[[\hat X_a,\hat X_b],A^a]\) is symmetric for \(\eta^{ab}\mathrm{Tr}(A_a^\dagger B_b)\), which forces a neutral eigenvector at any Jordan block, and it is symmetric for the Hilbert--Schmidt form when \(\mathrm{ad}[\hat X_0,\hat X_i]=0\). On a curved first-order background \([\hat X_a,\hat X_b]\) contains \(\mathsf R_{ab}{}^{cd}\mathcal S_{cd}\), and without vector fields commuting with all \(\mathbb D_a\) the first-order labels are absent. For random Hermitian \(3\times3\) matrices, not solutions, with \(\hat X_0\to\tau\hat X_0\), a complex pair forms at small \(\tau\) (about \(0.009\) at the seed used), with the two colliding eigenvectors tending to neutral (numerical). So the positivity of the pairing on shell and the absence of Jordan blocks on Hilbert--Schmidt kernels found above hold about commuting backgrounds and are not claimed for curved ones.

These statements are for Equation~\eqref{eq:Smatrix} with \(C_{\mu\nu\lambda}=0\), quadratic order about commuting matrices and the group cases named, the formal adjoints and kernels stated, with domains of the unbounded operators and the infrared volumes untreated. They do not derive that the matrix integral selects this class, or the frame and its non-degeneracy, so Assumptions~\ref{ass:frame-lock}, \ref{ass:nondeg}, \ref{ass:lorentzian} and~\ref{ass:local-lorentz} are not discharged. They are not a statement on the graviton two-point function of Section~\ref{sec:status-N} and do not discharge Assumption~\ref{ass:EFT-a2}: loops, masses of the labels, the dynamics of the diagonal elements, background torsion, sources, supersymmetry and a value or sign of \(G_{\mathrm{eff}}\) are not treated, and the frame of this class is not identified with the emergent metric of Chapter~\ref{ch:vierbein-emergence}. Problems~D, F and~N remain open and PROGRAMME-CONDITIONAL.

